\chapter*{Conclusiones}
\addcontentsline{toc}{chapter}{Conclusiones}

\section*{Resultados obtenidos en relación con la SAT}
\addcontentsline{toc}{section}{Resultados obtenidos en relación con la SAT}

El desarrollo del presente Proyecto Integrador ha culminado con el cumplimiento íntegro de los objetivos generales y específicos estipulados en la formulación de la investigación. Se logró la materialización técnica de un prototipo funcional de granja urbana vertical inteligente, capaz de monitorear y controlar de manera automatizada las variables críticas para el desarrollo vegetal (tales como la humedad del sustrato, la intensidad lumínica, la temperatura y la humedad relativa) mediante la aplicación práctica de principios de la Industria 4.0.  El sistema final satisface de manera rigurosa las especificaciones requeridas: se diseñó una estructura física modular y de disposición vertical para la optimización del espacio, se consolidó una arquitectura de procesamiento centralizada de alto rendimiento basada en una computadora de placa única (SBC), y se integraron subsistemas eficientes para la iluminación fotoperiódica, la irrigación localizada y la ventilación forzada. Asimismo, la validación empírica demostró la capacidad del prototipo para sostener el ciclo fenológico de las especies cultivadas, garantizando una gestión sustentable de los recursos hídricos mediante recirculación y certificando la viabilidad económica del sistema con un consumo energético acotado a plena carga.

\section*{Conclusión general}
\addcontentsline{toc}{section}{Conclusión general}

De manera global, el proyecto demuestra la viabilidad de diseñar e instrumentar sistemas de agricultura en ambiente controlado de bajo costo y alta precisión. La evolución de la plataforma, desde un prototipo preliminar de concepción estructural simplificada hasta la arquitectura modular definitiva, evidencia una pronunciada curva de aprendizaje que permitió identificar y resolver progresivamente múltiples barreras tecnológicas y fisicas.  A nivel mecánico e hídrico, las deficiencias tempranas de estanqueidad y drenaje, derivadas de la baja tensión superficial del material de impermeabilización, fueron sorteadas mediante el diseño y manufactura aditiva de acoples de desagüe roscados con sellos elastoméricos dobles. En el dominio electrónico y de instrumentación, la rápida degradación galvánica exhibida por los transductores resistivos obligó a una transición hacia sensores de humedad capacitivos, los cuales fueron impermeabilizados mediante carcasas impresas en 3D y sellado dieléctrico, asegurando la alta linealidad y estabilidad temporal de las mediciones edáficas.  Simultáneamente, los desafíos microclimáticos observados durante las primeras pruebas de cultivo (donde el estrés térmico indujo una floración prematura en las plantas) se mitigaron mediante una optimización en la lógica del software, la cual desplazó el fotoperíodo artificial hacia los ciclos nocturnos para favorecer la disipación térmica. Finalmente, la arquitectura informática maduró hacia un esquema de procesamiento tolerante a fallos, asegurando que la unidad central priorice el control local de los actuadores frente a eventuales pérdidas de conectividad con la plataforma telemétrica en la nube. Estas resoluciones integrales permitieron alcanzar la estabilidad operativa que garantizó el óptimo rendimiento biológico documentado en la fase final.

\section*{Aporte a la ingeniería o al campo de conocimiento}
\addcontentsline{toc}{section}{Aporte a la ingeniería o al campo de conocimiento}



\section*{Limitaciones y trabajos futuros}
\addcontentsline{toc}{section}{Limitaciones y trabajos futuros}

A pesar de los altos índices de confiabilidad y eficiencia documentados, la versión actual del prototipo presenta limitaciones físicas y operativas que definen líneas claras de continuidad para el proyecto.  En primer lugar, la adopción de sensores y transductores de bajo costo y alta disponibilidad comercial respondió a una estrategia metodológica para garantizar un prototipo accesible e imprevisto de bajos costos; sin embargo, esta elección evidenció restricciones en términos de precisión, repetibilidad y susceptibilidad a la deriva temporal en las mediciones. En segundo lugar, la interconexión de la red de sensado mediante cableado telefónico de cobre unifilar constituyó una vulnerabilidad mecánica importante; su rigidez y fragilidad ante manipulaciones o desplazamientos del módulo derivaron en fatiga del conductor y falsos contactos, degradando temporalmente la estabilidad del bus de datos. En tercer lugar, la etapa de potencia se resolvió reutilizando una fuente conmutada de factor de forma extendido que, si bien redujo los costos de implementación, introdujo un sobredimensionamiento volumétrico innecesario dentro del gabinete de control. Finalmente, la lógica algorítmica de actuación hídrica y microclimática se estructuró sobre esquemas de control discreto de dos posiciones (on/off) compensados mediante bandas de histéresis temporal, prescindiendo de modulación continua de caudal o potencia.  Para superar estas restricciones en futuras iteraciones del desarrollo, se contempla la actualización tecnológica integral del subsistema de percepción mediante la sustitución de los componentes actuales por sensores de especificación industrial o mayor rendimiento, garantizando un seguimiento microclimático de alta fidelidad. Asimismo, se propone el reemplazo del cableado unifilar por conductores multifilares apantallados con conectores de inserción rápida de grado industrial, e integrar un módulo de alimentación conmutado compacto y de perfil reducido que optimice el espacio físico de la envolvente. A nivel de procesamiento, la holgada capacidad computacional que ofrece la unidad central viabiliza el desarrollo de algoritmos de control predictivo, regulación de flujo continuo e incluso la incorporación de rutinas de visión computacional para el diagnóstico foliar no invasivo. Por último, la arquitectura informática desarrollada sienta las bases para la ampliación de los niveles físicos de cultivo y la adopción de protocolos de comunicación industrial avanzados, facilitando la escalabilidad del sistema hacia un modelo de producción comercial distribuido.